\section{引言：为什么需要句法结构}

当我们阅读或听到一句话时，表面上是一个线性的词序列（或线性的声音流），但人类大脑会自动地将其组织为具有层次结构的意义单元——判断哪些词修饰哪些词、哪些词组成一个短语。自然语言处理（NLP）系统同样需要理解句子的内部结构，才能正确地解释语言。

\begin{importantbox}{本讲核心内容}
\begin{enumerate}
    \item \textbf{句法结构的两种表示}：短语结构（Constituency/Phrase Structure）与依存结构（Dependency Structure）
    \item \textbf{依存语法的基本概念}：中心词（head）与依存词（dependent）、依存关系类型
    \item \textbf{句法歧义}：介词短语附着歧义、并列结构作用域歧义等
    \item \textbf{依存句法分析方法}：基于转移的分析（Transition-based Parsing）、基于图的分析（Graph-based Parsing）
    \item \textbf{神经网络依存分析器}：用分布式表示替代离散特征，实现更高精度与更快速度
\end{enumerate}
\end{importantbox}

本讲内容直接关联课程作业 Assignment 2 的第二部分——使用 PyTorch 构建一个神经依存句法分析器。

\begin{figure}[H]
\centering
\includegraphics[width=0.95\textwidth]{slides-images/slide-001.jpg}
\caption{CS224N Lecture 4 课程标题页\protect\footnotemark}
\end{figure}
\footnotetext{Slides 第1页。}

\subsection{本章小结}

句法结构是理解自然语言的基础。人类在理解语言时自动地解析句子的层次结构，NLP 系统也需要具备这种能力。本讲将从语言学概念出发，逐步引入计算方法。

%% ========== 短语结构语法 ==========

